\section{Architectural Guardrails for AI Agents}

The foundations chapter distinguished architecture erosion, caused by decisions that directly violate a governing architecture, from architecture drift, caused by decisions made without any awareness that a governing architecture exists at all --- and noted that this is a cumulative problem rather than a single-instance one: no individual decision needs to be wrong on its own terms for a codebase to gradually lose its coherence, since each decision may correctly satisfy the specific instruction it was given while still failing to align with a structural pattern established elsewhere. This distinction between erosion and drift matters for what kind of guardrail actually addresses the problem, because the two parts of the practical project differ precisely on whether a governing architecture was made explicit and enforceable in the first place.

Part 1's architectural problems are better characterized as drift than erosion: the missing service layer, the MainPage.js monolith, and the scattered reads of authentication state were not violations of a stated rule, because no rule existed yet to violate --- each decision was reasonable in isolation, and what was missing in every case was something persistent to check the decision against. \textbf{The guardrail this points to is not "write better individual instructions" but "establish an explicit, standing architecture before implementation begins," since drift specifically arises where no governing structure was ever articulated for a decision to be measured against.}

The Charter is Part 2's answer to this, and its rules are concrete rather than aspirational: a single application-factory pattern for the Flask application, exactly one blueprint file per resource domain, route handlers restricted to thin coordination rather than containing business logic, and a fixed return pattern for service functions rather than each one inventing its own error-handling convention. Several of these rules are traceable directly to a specific mistake already made in Part 1 --- the requirement that foreign keys to parent entities be non-nullable names the nullable \texttt{Collection.workspace\_id} foreign key from Part 1 as the violation not to repeat. \textbf{This is itself a best practice worth naming separately: a guardrail written as a direct response to a documented failure is more concrete and more enforceable than one derived abstractly, since it can be checked against a specific, known case rather than against a general principle open to interpretation.}

A rule's existence does not by itself guarantee conformance, and Part 2's structure addresses this by pairing the Charter with a role whose entire responsibility is checking adherence to it after code exists: the Post-Implementation Review Agent applies a fixed six-point checklist to every feature, covering categories --- untested code paths, excessive prop-drilling, unused or dead state, inconsistent naming --- that Part 1's process never checked at all, rather than only the two categories Part 1 caught occasionally and only once a problem had already spread. \textbf{The guardrail this establishes is that architectural rules need a standing, rule-based check paired with them to remain effective --- a rule that is stated once but never systematically verified affords no more protection against drift than having no rule at all, since an agent working under time or scope pressure has no structural reason to notice a quiet departure from it.}

